\section{Contributions}\label{section:conclusions-revisited}

The contributions of this thesis, supported by the results and discussion in Chapters~\ref{chapter:Results} and~\ref{chapter:Discussion}, are as follows: 

\begin{itemize}
    \item \textbf{An end-to-end, provenance-preserving extraction pipeline.} The pipeline turns histopathology PDFs from the PMC Open Access Subset into a structured clinical-knowledge base, in which each extracted rule, figure, and table is traceable to the source page, span, and coordinates it came from. Provenance is preserved by the data model from the beginning; it is structurally guaranteed. 

    \item \textbf{A document-extraction stage whose gains are measured against a baseline.} A hybrid Docling/TATR table detector, footnote expansion, and two-pass ghost-text filtering are combined in the document-extraction stage; this stage is evaluated under a rubric with four dimensions. The off-the-shelf Docling baseline isolates what the pipeline adds. For tables, capturing footnotes achieves the main improvement, raising strict table-F1 from $0.366$ to $0.838$. For figures, the main gain comes from three post-processing stages: decorative icon filtering, sub-figure merging, and custom caption matching (instead of Docling's default caption matcher); these post-processing stages raise strict figure-F1 from $0.447$ to $0.84$. These improvements are built on top of a base Docling parser, which already localizes tables and figures well. 

    \item \textbf{The agreement-based cascading (ABC) framework.} This thesis introduces and evaluates an agreement-based cascading framework with calibrated agreement scoring, accept/reject/escalate routing logic per chunk, single-voter and polarity hard-fail safety gates, and a pool of voters from multiple providers. The result is that, under the Opus silver-label metric, the cascade decisively beats cheaper single models, but is not cost-justified against the single strong model. Therefore, the cascade's main value does not come from a quality win against a single model, but rather from provider-independent voting and reusability of the framework.

    \item \textbf{A reproducible calibration and evaluation methodology.} The methodology comprises a calibration cluster that was selected by an integer linear program based on inter-paper relatedness, LLM silver labels, a matcher-based strict-F1 (with different matcher options), an explicit cost-quality frontier with different operating points (economy/knee/quality), and an offline-replay protocol that allows every reported number to be reproduced from cached artifacts, eliminating the need for additional API calls.

    \item \textbf{A shared, frozen NLI backbone.} A biomedical NLI model is used as-is, without any fine-tuning, and serves two roles in the pipeline: first, it checks whether extracted claims are supported by the evidence they cite; second, it classifies relations between two claims as supporting, contradicting, or unrelated. 

    \item \textbf{Frozen datasets used for empirical evaluation.} The document-extraction rubric set, the knowledge-extraction calibration and evaluation clusters, and the synthetic claim-pair dataset for relation classification make it possible to reproduce measurements and answer the research questions. 
\end{itemize}
